\section{De Bruijn Sequence}
\label{sec:cses-de-bruijn-sequence}

\problemstatement{For a given $n$, output a binary string that contains
\emph{every} binary string of length $n$ as a substring, and is as short as
possible (length $2^n+n-1$, treating the string as almost cyclic).}

\begin{patternbox}
A De Bruijn sequence is an \textbf{Eulerian path in the De Bruijn graph}:
nodes are the $2^{n-1}$ binary strings of length $n-1$, and each length-$n$
string is an edge from its $(n-1)$-prefix to its $(n-1)$-suffix. Every node
has equal in- and out-degree, so an Eulerian circuit exists.
\end{patternbox}

\subsection*{Edges are the length-n strings}

Build a graph whose nodes are all $(n-1)$-bit strings. For each $n$-bit
string, add an edge from its first $n-1$ bits to its last $n-1$ bits,
labelled by the last bit. Because appending $0$ or $1$ to any node gives a
valid edge, in-degree equals out-degree everywhere -- the Eulerian
condition. Run Hierholzer to traverse every edge once; concatenating the
starting node with each edge's appended bit produces a string in which every
length-$n$ substring appears exactly once.

\begin{ideabox}[Traverse Every Length-n String as an Edge]
Walking an Eulerian circuit uses each length-$n$ string (edge) exactly once,
and consecutive edges overlap in $n-1$ characters, so the emitted bits form
a compact sequence containing all $2^n$ patterns.
\end{ideabox}

\subsection*{Algorithm}

\begin{enumerate}
  \item Nodes $=$ integers $0\ldots2^{n-1}-1$; for each node $v$ and bit
        $b$, edge to $((v\ll1)\,|\,b)\bmod 2^{n-1}$ labelled $b$.
  \item Hierholzer from node $0$, recording edge labels.
  \item Output the start node's $n-1$ bits followed by the recorded labels.
\end{enumerate}

\subsection*{Dry run}

\dryrun{$n=2$: nodes $\{0,1\}$ (1-bit), edges for $00,01,10,11$.}

\begin{itemize}
  \item Eulerian circuit visits all four length-$2$ strings.
  \item A valid output is \texttt{00110}: substrings $00,01,11,10$ all
        appear.
\end{itemize}

\subsection*{C++ implementation}

\begin{lstlisting}
#include <bits/stdc++.h>
using namespace std;

int main() {
    int n;
    cin >> n;
    if (n == 1) { cout << "01\n"; return 0; }

    int nodes = 1 << (n - 1);
    int mask = nodes - 1;
    vector<int> ptr(nodes, 0);                    // next bit to try per node (0 then 1)

    // Hierholzer recording edge labels; stack holds {node, label used to enter}
    vector<int> labels;
    stack<pair<int,int>> st;
    st.push({0, -1});                              // start at node 0, no entering label
    while (!st.empty()) {
        auto [v, lab] = st.top();
        if (ptr[v] < 2) {                          // an unused out-edge (bit 0 then 1)
            int b = ptr[v]++;
            int to = ((v << 1) | b) & mask;
            st.push({to, b});
        } else {
            st.pop();
            if (lab != -1) labels.push_back(lab);
        }
    }
    reverse(labels.begin(), labels.end());

    string res(n - 1, '0');                        // start node = (n-1) zero bits
    for (int b : labels) res += char('0' + b);
    cout << res << "\n";
    return 0;
}
\end{lstlisting}

\begin{complexitybox}
\textbf{Time Complexity.} $O(2^n)$ -- one traversal of the $2^n$ edges.
\textbf{Space Complexity.} $O(2^n)$ for the recursion/circuit.
\end{complexitybox}

\begin{takeawaybox}
A De Bruijn sequence is an Eulerian circuit in the De Bruijn graph, where
length-$n$ strings are edges over $(n-1)$-length nodes. The construction is
the flagship application of Eulerian paths beyond literal edge-traversal
puzzles.
\end{takeawaybox}
